% ===========================================================================
%  Capítulo 1 — O Núcleo do Sistema
%  Prosa escrita à mão. As tabelas vêm de gen/.
% ===========================================================================
\chapter{O Núcleo do Sistema}

Este sistema abandona a ideia tradicional de Níveis de Personagem. Aqui, você não sobe de
nível automaticamente; o seu crescimento é orgânico, baseado no estudo, no treinamento e
no acúmulo de \rec{Pontos de Aprimoramento} (PA).

% ---------------------------------------------------------------------------
\section{Criação de Personagem e Atributos}

O sistema usa cinco atributos, e cada um é o bônus fixo que você soma nas rolagens de
d20:

\begin{description}
    \item[Força] poder físico, carga e ataques com armas pesadas.
    \item[Agilidade] reflexos, esquiva, furtividade e ataques precisos.
    \item[Vigor] saúde, resistência a venenos, clima e cansaço.
    \item[Intelecto] memória, conhecimento de magias e lógica.
    \item[Espírito] força de vontade, carisma, liderança e resistência mental.
\end{description}

\begin{nota}{Distribuindo seus Pontos Iniciais}
    Ao criar o personagem, você recebe \textbf{4 Pontos} para distribuir livremente entre
    os cinco atributos base. O valor máximo por atributo na criação é \textbf{4}.

    \medskip
    \textbf{Sistema de Defeitos:} você pode reduzir atributos pra ganhar pontos extras. As
    regras são estritas: apenas um atributo pode ficar em $-1$ (ganha 1 Ponto Extra) e
    apenas um em $-2$ (ganha 2 Pontos Extras). Nada desce abaixo de $-2$.

    \medskip
    Repare que o defeito \textbf{não cria pontos}: ele os empresta. Com os dois defeitos
    você distribui 7 pontos, mas dois atributos ficam em $-1$ e $-2$, então a \textbf{soma
    dos seus cinco atributos fecha em 4} de qualquer jeito. É essa soma --- não cada
    atributo isolado --- que a seção 2 usa pra cobrar PA por pontos comprados depois da
    criação. Consequência direta: \textbf{desfazer um defeito custa os mesmos 2 PA por
    ponto} que qualquer outro aumento.
\end{nota}

\begin{aviso}{Antes de largar o Vigor, leia isto}
    Vigor é o único atributo que não governa perícia nenhuma (seção 4), o que faz dele o
    alvo óbvio do Sistema de Defeitos. É de propósito que ele seja também o único cuja
    escala negativa é \textbf{desproporcional}: a Escala do Vigor (Cap. 4, §1) multiplica
    seus PV Máximos por \textbf{$\times$0,75} em $-1$ e por \textbf{$\times$0,40} em $-2$
    --- o primeiro ponto tira um quarto da sua vida, o segundo tira quase metade do que
    sobrou --- e ainda te dá Desvantagem em toda resistência de Vigor: veneno, doença,
    clima, fome, Exaustão e o Fio da Vida.

    \medskip
    Largar Força ou Intelecto custa um número na rolagem. Largar Vigor custa o personagem.
\end{aviso}

\begin{nota}{Os Dois Atributos do Mago}
    \textbf{Intelecto} --- a Precisão: define o quanto sua magia acerta e o quanto machuca.
    \textbf{Espírito} --- a Reserva: define quanto de mana o corpo consegue armazenar.

    \medskip
    \textbf{Escolas de magia não concedem PM.} A sua reserva inteira é \textbf{(o maior
    entre o seu Espírito e 4) $\times$ Maior Bônus de Rank de magia, $+$ 8} --- recalculada
    sempre que o seu maior patamar de magia sobe. Não há bônus extra por número de escolas
    abertas. Abrir oito escolas no 1º patamar não te dá mana nenhuma a mais; subir
    \emph{uma} escola até o Imperador multiplica tudo.

    \medskip
    Existem dois magos legítimos: o \textbf{cirurgião} (Intelecto alto --- poucos tiros,
    todos letais) e o \textbf{reator} (Espírito alto --- bombardeia o dia inteiro sem
    cansar). Roxy Migurdia é o primeiro. Rudeus Greyrat é o segundo.
\end{nota}

% ---------------------------------------------------------------------------
\section{Pontos de Aprimoramento (PA)}

A progressão ocorre quando o Mestre recompensa os jogadores com PA após sessões, missões
importantes ou arcos da história. Ao criar o personagem, você recebe \textbf{3 PA
iniciais}.

\begin{tabelalivro}{L{2.4cm} L{11cm}}{O que 1 PA compra fora das árvores}
\cabtabela{\textbf{Custo} & \textbf{O que você recebe}}
1 PA & 2 Perícias à sua escolha. \\
2 PA & $+$PV iguais a quatro vezes o seu maior Bônus de Rank (melhoria física permanente). \\
2 PA & $+$PM iguais ao dobro do seu maior Bônus de Rank de magia (melhoria mágica permanente). \\
2 PA & $+1$ ponto de Atributo Base permanente (teto 8) --- medido pela soma dos cinco, então desfazer um defeito custa o mesmo. \\
3 PA & Vantagem permanente em todos os Testes de Resistência de 1 Atributo à sua escolha (uma vez por atributo). \\
Variável & Magias, Técnicas e Talentos de Árvore --- o custo escala com o Rank (seção 3). \\
\end{tabelalivro}

\begin{nota}{Por que essas duas melhorias escalam com o Bônus de Rank}
    $+$PV e $+$PM da tabela acima não são um número fixo --- eles crescem junto com o seu
    maior Bônus de Rank (Principiante $+1$, até Imperador $+6$). Um Principiante gastando
    2 PA ganha pouco; um Imperador gastando os mesmos 2 PA ganha seis vezes mais. Isso
    evita que a compra vire golpe de sorte na criação e lixo de ficha no topo --- ela
    sempre pesa a mesma fração do que você já é.
\end{nota}

\begin{nota}{O Padrão das Reservas}
    Todo talento de árvore que compra reserva vale exatamente o mesmo, custe onde custar:
    \textbf{1 PA $=$ $+2$ PM por patamar seu naquela árvore}, ou \textbf{$+4$ PV por
    patamar seu naquela árvore}. No 1º patamar isso é $+2$ PM ou $+4$ PV, e é pouco; no 6º
    é $+12$ PM ou $+24$ PV, e vale a compra --- o mesmo talento cresce sozinho conforme
    você sobe naquela árvore.

    \medskip
    Isso existe pra que nenhuma escola dê mais vida de graça: escolas se diferenciam pela
    curva de progressão e pelas Maestrias, nunca por um talento genérico valer mais numa do
    que na outra.
\end{nota}

\begin{aviso}{Magia não tem preço fixo}
    Comprar Zero Absoluto não pode custar o mesmo que comprar Bola de Água. Toda aquisição
    dentro de uma Árvore de Progressão usa a tabela de custos por Rank da próxima seção.
\end{aviso}

% ---------------------------------------------------------------------------
\section{A Regra de Desbloqueio de Ranks}

O mundo é dividido em 7 Ranks de Maestria: Principiante, Intermediário, Avançado, Santo,
Rei, Imperador e Deus. Você só recebe permissão pra comprar o desbloqueio de um Rank
quando já possuir o número mínimo de \termo{conhecimentos} (magias e talentos comprados)
\emph{daquela mesma árvore}.

% GERADO por src/lib/livroTex.ts — não edite à mão; edite src/lib/types.ts.

\begin{tabelalivro}{L{2.6cm} C{2.2cm} C{2.6cm} C{2.2cm}}{Ranks de Maestria --- desbloqueio e exigência}
\cabtabela{\textbf{Rank} & \textbf{Bônus} & \textbf{Custo} & \textbf{Conhecimentos}}
Principiante & +1 & 1 PA & --- \\
Intermediário & +2 & 1 PA & 2 \\
Avançado & +3 & 2 PA & 4 \\
Santo & +4 & 2 PA & 6 \\
Rei & +5 & 3 PA & 8 \\
Imperador & +6 & 3 PA & 10 \\
Deus & — & Narrativa & — \\
\end{tabelalivro}

\begin{tabelalivro}{L{2.6cm} C{2.4cm} C{2.4cm} C{2.4cm}}{Custo em PA dentro de uma árvore}
\cabtabela{\textbf{Rank} & \textbf{Comum} & \textbf{Assinatura \assinatura{}} & \textbf{Talento}}
Principiante & 1 PA & 2 PA & 1 PA \\
Intermediário & 1 PA & 2 PA & 1 PA \\
Avançado & 2 PA & 3 PA & 2 PA \\
Santo & 3 PA & 4 PA & 3 PA \\
Rei & 4 PA & 5 PA & 4 PA \\
Imperador & 5 PA & 6 PA & 4 PA \\
\end{tabelalivro}

\begin{tabelalivro}{L{2.6cm} C{2.8cm} C{2.8cm}}{Custo em PA nas árvores de Utilidade (tabela própria, mais barata)}
\cabtabela{\textbf{Rank} & \textbf{Talento} & \textbf{Assinatura \assinatura{}}}
Principiante & 1 PA & 2 PA \\
Intermediário & 1 PA & 2 PA \\
Avançado & 2 PA & 3 PA \\
Santo & 2 PA & 3 PA \\
Rei & 3 PA & 4 PA \\
Imperador & 3 PA & 4 PA \\
\end{tabelalivro}

\begin{tabelalivro}{L{2.6cm} C{2.2cm} C{2.4cm} C{3.2cm}}{Tempo de conjuração por rank}
\cabtabela{\textbf{Rank da magia} & \textbf{Padrão} & \textbf{Encurtada} & \textbf{Silenciosa}}
Principiante & 2 Ações & 1 Ação(ões) & 1 Ação(ões) (1ª grátis) \\
Intermediário & 2 Ações & 1 Ação(ões) & 1 Ação(ões) \\
Avançado & 3 Ações & 2 Ação(ões) & 1 Ação(ões) \\
Santo & 4 Ações & 3 Ação(ões) & 2 Ação(ões) \\
Rei & 5 Ações & 4 Ação(ões) & 3 Ação(ões) \\
Imperador & 6 Ações & Impossível & 4 Ação(ões) \\
\end{tabelalivro}

\begin{tabelalivro}{C{1.6cm} L{3.4cm} C{2.2cm} L{5.4cm}}{A Escala do Vigor}
\cabtabela{\textbf{Vigor} & \textbf{Nome} & \textbf{Fator de PV} & \textbf{O que significa}}
-2 & Corpo Quebrado & $\times$0,40 & 60\% da vida de um corpo comum. \\
-1 & Constituição Frágil & $\times$0,75 & 75\% da vida de um corpo comum. \\
+0 & Corpo Comum & $\times$1,00 & A referência. \\
+1 & Saudável & $\times$1,20 & +20\% sobre o corpo comum. \\
+2 & Robusto & $\times$1,40 & +40\% sobre o corpo comum. \\
+3 & Endurecido & $\times$1,60 & +60\% sobre o corpo comum. \\
+4 & Inquebrável & $\times$1,80 & +80\% sobre o corpo comum. \\
+5 & Sobre-humano & $\times$2,00 & +100\% sobre o corpo comum. \\
+6 & Monstruoso & $\times$2,20 & +120\% sobre o corpo comum. \\
+7 & Lendário & $\times$2,40 & +140\% sobre o corpo comum. \\
+8 & Divino & $\times$2,60 & +160\% sobre o corpo comum. \\
\end{tabelalivro}

\begin{tabelalivro}{C{1.4cm} L{11cm}}{A Escada de Dados de Arma (Cap. 3)}
\cabtabela{\textbf{Degrau} & \textbf{Dado}}
0 & \dado{d4} \\
1 & \dado{d6} \\
2 & \dado{d8} \\
3 & \dado{d10} \\
4 & \dado{d12} \\
5 & \dado{2d8} \\
6 & \dado{2d10} \\
7 & \dado{2d12} \\
8 & \dado{3d10} \\
9 & \dado{3d12} \\
10 & \dado{4d10} \\
11 & \dado{4d12} \\
12 & \dado{5d10} \\
13 & \dado{5d12} \\
\end{tabelalivro}

\begin{tabelalivro}{L{6cm} C{2cm}}{Dados Base de arma}
\cabtabela{\textbf{Arma} & \textbf{Dado}}
Adaga / Punhal & \dado{d4} \\
Funda / Dardo & \dado{d4} \\
Chicote & \dado{d4} \\
Espada Curta & \dado{d6} \\
Objeto Improvisado & \dado{d6} \\
Arco Curto & \dado{d6} \\
Rapieira & \dado{d6} \\
Espada Longa & \dado{d8} \\
Machado de Batalha & \dado{d8} \\
Arco Longo & \dado{d8} \\
Foice de Guerra & \dado{d8} \\
Espadão / Montante & \dado{d10} \\
Martelo de Guerra & \dado{d10} \\
Alabarda / Lança & \dado{d10} \\
Besta & \dado{d10} \\
\end{tabelalivro}

\begin{tabelalivro}{C{1.8cm} L{4.2cm} L{7.6cm}}{Perícias de Árvore --- o que cada uma das 18 ensina (Cap. 1, §4)}
\cabtabela{\textbf{Pilar} & \textbf{Árvore} & \textbf{Ensina, se for a sua Árvore Inicial}}
Magia & Magia de Fogo & Arcanismo --- mais 1 à sua escolha entre Ofícios, Intimidação e Atletismo. \\
Magia & Magia de Água & Arcanismo --- mais 1 à sua escolha entre Medicina, Natureza e Percepção. \\
Magia & Magia de Vento & Arcanismo --- mais 1 à sua escolha entre Acrobacia, Percepção e Sobrevivência. \\
Magia & Magia de Terra & Arcanismo --- mais 1 à sua escolha entre Ofícios, Atletismo e Natureza. \\
Magia & Magia de Cura & Medicina --- mais 1 à sua escolha entre Religião, Intuição e Persuasão. \\
Magia & Magia de Desintoxicação & Medicina --- mais 1 à sua escolha entre Natureza, Ofícios e Investigação. \\
Magia & Barreira e Proteção & Arcanismo --- mais 1 à sua escolha entre Religião, Percepção e Intuição. \\
Magia & Espíritos e Feras & Arcanismo --- mais 1 à sua escolha entre Natureza, Lidar com Animais e Religião. \\
Corpo & Estilo Deus da Espada & Acrobacia e Intimidação. \\
Corpo & Estilo Deus da Água & Intuição e Percepção. \\
Corpo & Estilo Deus do Norte & Sobrevivência e Enganação. \\
Corpo & Lutador & Atletismo e Intimidação. \\
Corpo & Cavalaria e Escudos & Atletismo e Percepção. \\
Corpo & Estilo Vendaval & Acrobacia e Percepção. \\
Corpo & Arquearia & Percepção e Sobrevivência. \\
Utilidade & Furtividade e Armadilhas & Furtividade e Percepção --- mais 1 à sua escolha entre Ladinagem, Acrobacia e Enganação. Exceção única do livro: se esta NÃO for a sua Árvore Inicial, a Maestria de 1º patamar ainda ensina Furtividade e Percepção. \\
Utilidade & Bardo e Interação & Atuação e Persuasão --- mais 1 à sua escolha entre Intuição e História. \\
Utilidade & Navegação e Liderança & Sobrevivência e Investigação --- mais 1 à sua escolha entre Natureza e Percepção. \\
\end{tabelalivro}


\begin{nota}{Magia Assinatura \assinatura}
    Dentro de cada Rank existe uma magia que define aquele patamar --- a que os magos
    daquele nível são reconhecidos por saber. Ela custa $+1$ PA a mais que uma magia comum
    do mesmo rank, e é marcada com \assinatura{} em todas as listas do Capítulo 3.
\end{nota}

\begin{nota}{Maestrias não contam}
    Maestrias (as passivas automáticas ganhas de graça ao desbloquear um Rank) não contam
    como conhecimento. Apenas magias e talentos efetivamente comprados com PA contam pra
    tabela acima.
\end{nota}

\begin{aviso}{E o Rank Deus}
    O patamar Divino não possui custo mecânico de PA. Como habilidades divinas beiram a
    onipotência e reescrevem as leis da realidade, este Rank só pode ser alcançado através
    de intenso Roleplay e eventos lendários na narrativa, ditados inteiramente pela
    história e pelo Mestre.
\end{aviso}

% ---------------------------------------------------------------------------
\section{O Sistema de Testes e Perícias}

Sempre que um jogador tentar uma ação com chance de falha, ele rola \textbf{1d20 $+$ o
Atributo correspondente}.

\begin{itemize}
    \item \textbf{Compra Direta:} 1 PA adquire 2 Perícias simultaneamente.
    \item \textbf{Perícias de Árvore:} certas árvores e talentos concedem perícias
          específicas como bônus por treinar aquele estilo.
\end{itemize}

\begin{nota}{Vantagem por Perícia}
    Sempre que você realizar uma ação e possuir uma Perícia que se encaixe perfeitamente
    na situação, você recebe \termo{Vantagem}: role 2d20, escolha o maior resultado, e só
    então some o bônus do Atributo Base.

    \medskip
    \textbf{Vantagem Absoluta:} alguns efeitos concedem Vantagem Absoluta --- role 3d20 e
    escolha o maior. Desvantagem Absoluta funciona igual, mas escolhendo o menor.
\end{nota}

\subsection{Lista Mestre de Perícias}

Vinte perícias, cada uma sob o atributo que a testa. Lista fechada --- sem encaixe
perfeito, o teste é só o Atributo puro. \textbf{Vigor não governa nenhuma perícia}: ele já
é a reserva de PV e a resistência a veneno, clima e cansaço do Capítulo 4.

% GERADO — fonte: src/data/skills.ts
\begin{tabelalivro}{L{2.2cm} L{3.2cm} L{8.4cm}}{Lista Mestre de Perícias}
\cabtabela{\textbf{Atributo} & \textbf{Perícia} & \textbf{Cobre}}
Força & Atletismo & Escalar, nadar, arrombar, forçar, saltar. \\
Agilidade & Acrobacia & Equilíbrio e esquiva acrobática. \\
Agilidade & Furtividade & Esconder-se e mover-se sem ser notado. \\
Agilidade & Ladinagem & Agilidade manual --- bater carteira, soltar algemas, plantar um item. \\
Intelecto & Arcanismo & Teoria mágica e itens encantados. \\
Intelecto & História & Fatos do passado e linhagens. \\
Intelecto & Investigação & Deduzir pistas e ligar evidências. \\
Intelecto & Medicina & Primeiros socorros e diagnóstico físico. \\
Intelecto & Natureza & Fauna, flora e clima. \\
Intelecto & Ofícios & Um ofício manual específico, escolhido ao adquirir (Forja, Culinária, Alquimia, Carpintaria...). \\
Intelecto & Religião & Doutrina, templos e o Continente Divino. \\
Espírito & Atuação & Performance, música, oratória de palco. \\
Espírito & Enganação & Mentir e disfarçar intenção. \\
Espírito & Intimidação & Impor medo. \\
Espírito & Intuição & Perceber mentira e prever intenção. \\
Espírito & Lábia & Convencer rápido, pechinchar, tagarelar. \\
Espírito & Lidar com Animais & Acalmar e comandar animais. \\
Espírito & Percepção & Notar detalhes com os sentidos. \\
Espírito & Persuasão & Convencer com argumento sincero. \\
Espírito & Sobrevivência & Rastrear, orientar-se e sobreviver no ermo. \\
\end{tabelalivro}


\begin{nota}{Sobre nomes parecidos}
    Persuasão é o argumento sincero. Lábia é o oposto: rapidez de fala, pechincha de
    mercado, o papo que convence sem precisar ser verdadeiro. Um Antecedente ou traço que
    mencione \enquote{Diplomacia} se refere a Persuasão --- o livro usa um nome só.
\end{nota}

\subsection{Proficiências: Armas e Armaduras}

Toda árvore do Corpo já concede proficiência com o que ela usa. O padrão abaixo vale pra
quando nenhuma árvore ainda cobriu aquele equipamento.

\begin{itemize}
    \item \textbf{Armas simples} (Dado Base até d6): todo personagem é proficiente.
    \item \textbf{Armas marciais} (Dado Base d8$+$): exigem proficiência do 1º patamar de
          uma árvore do Corpo, ou de um talento específico.
    \item \textbf{Armadura leve:} todo personagem é proficiente.
    \item \textbf{Armadura média/pesada e escudos:} exigem proficiência específica.
\end{itemize}

\begin{aviso}{Penalidade de Não-Proficiência}
    \textbf{Arma sem proficiência:} Desvantagem no teste de acerto. O dano continua normal
    --- a Escada de Dados nunca reduz.

    \medskip
    \textbf{Armadura sem proficiência:} Desvantagem em Furtividade e Acrobacia, e
    Deslocamento $-3$m enquanto vestida.
\end{aviso}

\subsection{Equipamento Inicial e a Árvore Inicial}

Ninguém começa do zero. Ao criar seu personagem, desbloqueie o 1º patamar de pelo menos
uma árvore --- sua \termo{Árvore Inicial} --- com parte dos seus 3 PA iniciais. Ela decide
o kit abaixo, recebido de graça e \emph{além} do dinheiro do Antecedente: as duas coisas
não competem entre si.

% GERADO — fonte: src/data/startingKits.ts
\begin{tabelalivro}{L{4.2cm} L{9.6cm}}{Equipamento Inicial por Árvore Inicial}
\cabtabela{\textbf{Árvore Inicial} & \textbf{Kit}}
Estilo Deus da Espada, Estilo Deus da Água, Estilo Deus do Norte & Espada Longa · Roupas de viagem resistentes \\
Lutador & Machado de Batalha \\
Cavalaria e Escudos & Escudo · Espada Curta · Armadura Média \\
Arquearia & Arco Curto (20 flechas inclusas.) · Adaga reserva \\
Magia de Fogo, Magia de Água, Magia de Vento, Magia de Terra & Cajado / Foco Arcano (Conta como objeto improvisado corpo a corpo.) · Grimório básico · Roupas de viajante \\
Magia de Cura, Magia de Desintoxicação, Barreira e Proteção & Kit de cura (Bandagens, ervas, um frasco vazio.) · Símbolo do templo de origem · Cajado leve \\
Espíritos e Feras & Giz e tinta ritual (Material pra três círculos.) · Adaga \\
Furtividade e Armadilhas & Adaga · Adaga reserva · Kit de arrombamento · Roupas escuras \\
Bardo e Interação & Instrumento musical (Escolha o instrumento.) · Adaga \\
Navegação e Liderança & Kit de sobrevivência (Corda, mapa em branco, uma semana de provisões.) · Arma simples (Escolha qual.) \\
\end{tabelalivro}


\begin{nota}{O kit não é a build}
    O kit inicial existe só pra ninguém chegar na primeira cena de mãos vazias. Ele não
    substitui comprar magias, técnicas ou talentos com PA, e pode ser vendido, trocado ou
    ignorado como qualquer outro item da mochila.
\end{nota}

% ---------------------------------------------------------------------------
\section{Raças do Mundo de Seis Faces}

O mundo é habitado por raças com fisiologias e culturas vastamente diferentes. Escolha sua
linhagem pra determinar traços genéticos e mecânicos.

% GERADO — fonte: src/data/races.ts
\subsection{Humano (Jinzoku)}
A raça dominante do mundo. Físico relativamente fraco e vida curta (70-100 anos), mas altíssima inteligência e versatilidade.

\begin{itemize}
  \item Adaptabilidade: 2 Perícias extras à escolha e +1 em UM atributo à sua escolha, permanente --- nenhuma outra raça deixa você decidir onde o sangue pesa.
  \item Determinação Humana: uma vez por sessão, repita um teste de Atributo (não de Perícia, não de dano) que tenha acabado de falhar e use o novo resultado --- humanos vivem menos que qualquer raça deste livro e aprenderam a não desperdiçar a única tentativa que têm.
\end{itemize}

\subsection{Elfo (Erufu)}
Habitantes da Grande Floresta. Corpos esguios, orelhas longas, fertilidade baixa e vida longuíssima.

\begin{itemize}
  \item Sentido da Floresta: Vantagem em Percepção auditiva e em Sobrevivência para navegação.
  \item +1 em Agilidade, permanente, e PM Máximos iguais ao DOBRO do seu Maior Bônus de Rank de magia (+2 no Principiante, +12 no Imperador) --- séculos de convivência com a mana da Grande Floresta. Sem nenhuma escola de magia aberta, este bônus é 0.
  \item Sangue Longevo: Vantagem em testes de resistência de Vigor contra veneno e doença (Cap. 4, §7) --- séculos de vida ensinam o corpo a esperar o pior.
\end{itemize}

\subsection{Anão (Dowaafu)}
Artesãos e ferreiros inatos da Cordilheira do Dragão Azul. Vivem várias centenas de anos, baixa estatura, alta resistência ao álcool.

\begin{itemize}
  \item Sangue da Forja: magias de Terra e Fogo custam 1 PM a menos para conjurar (mínimo 1). Não pode aprender magias de Água ou Vento.
  \item +1 em Vigor e +10 PV Máximos, permanentes --- o corpo mais denso do livro.
  \item Fígado de Pedra: imune a ficar Embriagado e tem Vantagem em testes de resistência de Vigor contra Exaustão por privação (Cap. 4, seção 8).
\end{itemize}

\subsection{Povo Pequeno / Hobbit (Hobitto)}
Vivem na Grande Floresta e em cidades como Millishion. Estatura e aparência de criança humana por toda a vida.

\begin{itemize}
  \item Deslocamento base reduzido: 7,5m.
  \item Aparência Enganosa: Vantagem em Enganação e Furtividade.
  \item Pequeno Demais pra Atrapalhar: você pode ocupar o mesmo espaço de outra criatura, desde que ela permita --- passar por baixo, subir no ombro, se enfiar atrás das pernas dela. Não concede Cobertura automática nem impede que você seja alvo; só deixa vocês dois no mesmo quadrado.
  \item Sombra Absoluta (opcional, 3 PA): transforme a Vantagem racial acima em Vantagem Absoluta (Cap. 1, §4). É uma compra, não um bônus grátis --- e é a única melhoria racial comprável do livro.
  \item +1 em Agilidade, permanente.
  \item Sorte do Povo Pequeno: uma vez por Descanso Longo, transforme uma Falha Crítica (1 Natural) sua em um resultado normal --- o dado ainda rola, mas o desastre automático não acontece.
\end{itemize}

\subsection{Raça Fera (Juuzoku)}
Habitantes da Grande Floresta com traços de mamíferos. Fisicamente superiores aos humanos, vida similar.

\begin{itemize}
  \item Sentidos Selvagens: Vantagem para rastrear pelo olfato; Desvantagem em resistência a fumaça/odores fortes.
  \item Magia Inerente --- HOWLING (2 PM, 1 Ação): você nasce sabendo, sem gastar PA e sem precisar de escola aberta. Ao conjurar, escolha UM dos dois modos.
  \item Howling · Grito de Guerra (ataque sônico): cone de 9 metros. Cada criatura na área faz teste de resistência de Vigor contra CD 8 + Espírito + seu Maior Bônus de Rank. Falha: sofre 1d6 de dano sônico por ponto do seu Maior Bônus de Rank (1d6 no Principiante, 6d6 no Imperador) e fica Atordoada até o fim do próximo turno dela. Sucesso: metade do dano e nada mais. Uma vez por combate --- depois do primeiro uivo, ninguém mais é pego de surpresa.
  \item Howling · Eco de Caça (ecolocalização): o uivo volta e desenha o que tocou. Por 1 minuto você sabe a posição exata de toda criatura a até 30 metros, mesmo no escuro total, mesmo sob invisibilidade mágica, mesmo através de porta, mato ou parede fina. Você sabe ONDE, nunca O QUÊ: tamanho aproximado e posição, não identidade nem intenção. Pedra maciça, chumbo e qualquer barreira mágica bloqueiam o eco. Sem limite de usos.
  \item +1 em Força, permanente.
  \item Instinto de Caçada: Vantagem em Iniciativa contra qualquer criatura que você tenha farejado, rastreado ou observado antes do combate começar.
\end{itemize}

\subsection{Raça Celestial (Tenzoku)}
Habitantes do Continente Divino. Vivem centenas de anos e possuem asas.

\begin{itemize}
  \item Deslocamento de Voo igual ao de caminhada --- só sem armadura média ou pesada e sem carregar mais da metade do seu limite de carga: asas não erguem aço.
  \item +1 em Espírito, permanente.
  \item Sangue do Continente Divino: Vantagem em testes de resistência de Espírito contra Medo e contra qualquer efeito de origem divina.
\end{itemize}

\subsection{Raça do Oceano (Kaizoku)}
Governantes do Mar de Ringus.

\begin{itemize}
  \item Respira debaixo d'água.
  \item Ignora penalidades de terreno difícil aquático.
  \item +1 em Vigor, permanente.
  \item Pressão das Profundezas: Resistência a dano contundente vindo de água em movimento (correnteza, magia de Água que usa força bruta, tsunami de cerco).
\end{itemize}

\subsection{Migurd}
Humanoides de cabelos e olhos azuis, \textasciitilde{}200 anos de vida, aparência de adolescente até os 150 anos.

\begin{itemize}
  \item +1 em Intelecto, permanente, e PM Máximos iguais ao TRIPLO do seu Maior Bônus de Rank de magia (+3 no Principiante, +18 no Imperador) --- a maior reserva racial do livro, e a única que não decai. Sem nenhuma escola de magia aberta, este bônus é 0.
  \item Telepatia curta com outros Migurds ou seres com telepatia.
  \item Mente Fechada: Vantagem em testes de resistência de Espírito contra qualquer efeito que leia, controle ou confunda a mente --- quem nasce falando por telepatia aprende a trancar a própria porta antes de aprender a andar.
\end{itemize}

\subsection{Superd}
Pele pálida, cabelos verdes, cauda bifurcada que vira lança tridente.

\begin{itemize}
  \item Previsão de Movimento (1 Ação): escolha uma criatura a até 18m que você possa ver e leia o fluxo de mana dela por 1 minuto --- você enxerga o golpe antes de ele sair. Enquanto durar: os ataques dela contra você têm Desvantagem, você tem Vantagem nos testes de resistência contra as habilidades dela, e ela nunca te pega Surpreso.
  \item Previsão de Movimento --- limites: uma leitura por vez (trocar de alvo custa outra Ação), e não funciona contra o que não move mana: construto inerte, armadilha mecânica, uma pedra caindo. Ler o fluxo não é ver o futuro; é ver a intenção antes de ela virar movimento.
  \item Sofre Desvantagem em interações sociais com humanos comuns (preconceito antigo, Cap. 1) --- mas Vantagem Absoluta em Intuição para perceber a intenção real de quem esconde algo, porque o Terceiro Olho não mente.
\end{itemize}

\subsection{Ogro (Onizoku)}
Extremamente altos e musculosos, machos chegam a 3 metros de altura.

\begin{itemize}
  \item Brutamontes: Vantagem em testes de Força bruta.
  \item Limite de carga dobrado.
  \item +2 em Força, permanente --- o maior bônus de atributo de qualquer raça do livro.
\end{itemize}

\subsection{Demônio Imortal}
Descendentes do Primeiro Deus Demônio. Pele negra azeviche, seis braços (machos).

\begin{itemize}
  \item Regeneração Profunda: regenera PV iguais ao seu Maior Bônus de Rank (Cap. 1, §7) no início do seu turno, desde que esteja com mais de 0 PV.
  \item +8 PV Máximos, permanentes.
  \item Descendência Divina: Vantagem em testes de resistência de Vigor contra veneno e doença (Cap. 4, §7).
\end{itemize}

\subsection{Raça Dragão (Ryuzoku)}
Raça mítica (requer aprovação do Mestre). Fisicamente a mais poderosa da existência, pode viver mais de 100.000 anos.

\begin{itemize}
  \item Escamas Dracônicas: +3 na CA (permanente, empilha com armadura) e Resistência a dano cortante e perfurante não-mágico.
  \item Garras e Presas: seus ataques desarmados usam Dado Base d10, contam como arma marcial mágica e sobem na Escada de Dados (Cap. 3) junto com o seu maior patamar do Corpo --- um Dragão desarmado nunca está desarmado.
  \item Sopro Dracônico (1 Ação, 1 vez por Descanso Curto): cone de 12 metros do elemento que você escolheu ao criar o personagem. Dano igual a 1d10 por ponto do seu Maior Bônus de Rank (1d10 no Principiante, 6d10 no Imperador). Teste de resistência de Agilidade contra CD 8 + Vigor + Maior Bônus de Rank para metade do dano.
  \item Asas: Deslocamento de Voo igual ao dobro do seu Deslocamento de caminhada, sem restrição de armadura ou carga --- as asas de um Ryuzoku erguem aço sem esforço.
  \item +2 em Força e +1 em Vigor, permanentes, e IMUNIDADE (não Resistência) a um elemento à escolha: Fogo, Gelo ou Eletricidade. É o mesmo elemento do seu Sopro.
  \item Cem Mil Anos: você não envelhece de forma perceptível, é imune a doença comum, e tem Vantagem em testes de resistência de Espírito contra qualquer efeito de Medo ou de controle mental.
  \item O Preço do Sangue: Vantagem em Intimidação, mas Desvantagem Absoluta em Persuasão, Lábia ou Diplomacia --- nada que já foi um deus finge ser gente comum de verdade. É o único preço que a raça cobra, e ele é permanente.
\end{itemize}



% ---------------------------------------------------------------------------
\section{O Destino e a Infância (Antecedentes)}

O que você fez nos seus primeiros 10 anos de vida define a fundação do seu corpo, sua mana
e seu lugar no mundo. Durante a criação, role 1d100 (ou escolha com o Mestre) pra descobrir
sua origem e seu dinheiro inicial em Peças de Ouro (PO).

Na sociedade humana, as anomalias de mana que causam poderes são chamadas de \termo{Miko}
(Criança Abençoada) se o poder for útil, ou \termo{Noroi-ko} (Criança Amaldiçoada) se for
prejudicial. Existem apenas cerca de 10 Mikos em todo o mundo simultaneamente.

% GERADO — fonte: src/data/backgrounds.ts
\begin{tabelalivro}{C{1.3cm} L{3.1cm} L{7.6cm} C{1.7cm}}{O Destino e a Infância (1d100)}
\cabtabela{\textbf{d100} & \textbf{Antecedente} & \textbf{Efeito} & \textbf{Dinheiro}}
01--15 & Plebeu / Trabalhador Rural & 2 Perícias à escolha · 2 Perícias ligadas a trabalhos mundanos (Ofícios, Culinária, Lidar com Animais ou Natureza). · +1 em Vigor, permanente --- corpo calejado por uma infância inteira de trabalho braçal. & 2d4 PO \\
16--25 & Órfão das Ruas & 2 Perícias à escolha · 2 Perícias de sobrevivência urbana (Furtividade, Ladinagem, Enganação ou Acrobacia). · +1 em Agilidade, permanente. · Instinto de Rua: Vantagem em Percepção para notar armadilhas, emboscadas, bolsos alheios e vigias --- ninguém sobrevive na rua sem aprender a ler uma esquina. & 1d4 PO \\
26--35 & Criança Selvagem & Perícia: Sobrevivência · +6 PV Máximos e +1 em Vigor, permanentes. · Rola Sobrevivência com Vantagem para achar comida, água e abrigo. & 0 PO \\
36--45 & Aprendiz de Mercador & Perícia: Intuição · Perícia: Lábia · +1 em Intelecto, permanente. · Sexto Sentido Comercial: sempre sabe o preço justo de mercado de qualquer item comum, e tem Vantagem em testes pra perceber quando alguém está blefando numa negociação. & 4d4+10 PO \\
46--55 & Treino Precoce / Escudeiro & Perícia: Atletismo · Vantagem em todas as rolagens de Iniciativa. · +1 em Força, permanente --- anos carregando a armadura de outra pessoa antes de vestir a própria. & 2d4+2 PO \\
56--65 & Acólito / Filho do Templo & Perícia: Religião · Perícia: Medicina Básica · +1 em Espírito e +4 PM Máximos, permanentes --- a educação religiosa desperta uma afinidade latente com mana que nunca mais desaparece. & 2d4 PO \\
66--72 & Sangue Nobre & Perícia: Persuasão · Perícia: História · Vantagem em testes sociais ao lidar com autoridades. · +1 em Espírito, permanente --- comandar serviçais desde criança ensina presença antes de ensinar humildade. & 6d4+20 PO \\
73--80 & Estudioso Precoce (Expansão) & Perícia: Arcanismo · +8 PM Máximos e +1 em Intelecto, permanentes. & 2d4 PO \\
81--86 & Sobrevivente / Ex-Escravo & Vantagem em resistência de Espírito (contra Medo) e de Vigor (contra Exaustão). · +1 em Vigor e +1 em Espírito, permanentes --- o corpo que sobreviveu ao pior já não se assusta com o segundo pior, e a vontade que o carregou até aqui não some junto com as cicatrizes. & 0 PO \\
87--92 & Fator Laplace / Linhagem Antiga & +2 em Espírito, +8 PM Máximos e +6 PV Máximos, permanentes --- seu corpo nasceu acostumado a segurar mais mana e mais dor do que deveria ser possível. · Conjuração Silenciosa desde o nascimento (Cap. 2, seção 2): você manipula mana sem palavra alguma --- ninguém te ensinou, você nunca soube fazer diferente. · Vantagem em testes de resistência de Espírito contra Medo, Amedrontado e qualquer efeito que tente controlar sua mente: o que quer que exista na sua linhagem, não se deixa comandar. · Desvantagem em Persuasão com desconhecidos --- pessoas comuns sentem, mesmo sem saber o porquê, que algo em você quer distância. & 1d4 PO \\
93--96 & Miko (Abençoada/Amaldiçoada) & +1 em Espírito, permanente --- a anomalia de mana que produz a mutação não fica só no poder: ela engrossa a alma que o carrega. · Role 1d8 na Tabela de Miko para definir a mutação mágica. & 2d4 PO \\
97--98 & Olho Místico Inato & +1 em Intelecto e +6 PM Máximos, permanentes --- o olho não vem sozinho: o corpo que o carrega já nasce com a mana de sobra que ele consome. · Role 1d10 na Tabela de Olhos Místicos para definir o Magan. & 2d4 PO \\
99--100 & Gênio (Conjuração Silenciosa) & Conjuração Silenciosa (Cap. 2, seção 2) sem sofrer a redução de dano nem a redução de área --- só a limitação de forma continua valendo, e mesmo essa você pode escolher mudar como qualquer outro conjurador silencioso. · +1 em Intelecto e +1 em Espírito, permanentes --- precisão e reserva ao mesmo tempo, que é o que fez de você um gênio antes de qualquer aula. & 2d4 PO \\
\end{tabelalivro}

\begin{tabelalivro}{C{1.1cm} L{3.6cm} L{9cm}}{Miko e Amaldiçoados (1d8)}
\cabtabela{\textbf{1d8} & \textbf{Tipo} & \textbf{Efeito}}
1 & Força Sobre-humana (Zanoba) & Abençoada: +2 em Força, permanente, e ataques desarmados causam 1d8 + Força (letal). Maldição: Desvantagem em testes de Vigor (sem resistência física). \\
2 & Leitura de Memórias & Abençoada: 1 Ação de toque lê memórias superficiais e intenções (Vantagem em Intuição/Interrogatório). Maldição: custa 3 PM por uso. \\
3 & Rebobinar o Tempo & Abençoada: 1x/semana rebobina o estado de um objeto inanimado em até 24h. Maldição: drena 50\% do PM Máximo atual; não funciona em criaturas vivas. \\
4 & Telepatia & Abençoada: lê pensamentos superficiais e fala telepaticamente num raio de 18m. Maldição: fisicamente muda; Desvantagem em Iniciativa. \\
5 & Miko da Confiança Absoluta & Abençoada: Vantagem Absoluta em Persuasão e Lábia. Maldição: falha automaticamente em Intuição para perceber mentiras. \\
6 & Maldição do Esquecimento & Abençoada: Vantagem Absoluta em Furtividade (presença nula). Maldição: quase ninguém lembra do seu rosto/nome/existência 10min após sair do campo de visão. \\
7 & Maldição do Acúmulo & Abençoada: +10 PM Máximos (reator infinito). Maldição: exige 'liberação' semanal; falhar aplica 1 nível de Exaustão por dia até a morte. \\
8 & Maldição do Ódio & Abençoada: +2 na CA e +10 PV Máximos (aura primordial). Maldição: todo ser que sinta mana sofre ódio instintivo e paranóico ao te ver. \\
\end{tabelalivro}

\begin{tabelalivro}{C{1.1cm} L{3.6cm} L{9cm}}{Olhos Demoníacos / Místicos (1d10)}
\cabtabela{\textbf{1d10} & \textbf{Olho} & \textbf{Mecânica}}
1 & Olho da Previsão & 1 Ação pra ativar, 3 PM/turno: vê 2s no futuro (Vantagem em ataques, oponentes têm Desvantagem contra você). Mais de 3 turnos seguidos causa Tontura por 1h. \\
2 & Olho do Poder Mágico & 1 Ação, 2 PM/cena: vê fluxo de mana, invisíveis mágicos, identifica itens mágicos e nível de perigo. \\
3 & Olho da Clarividência & 1 Ação, 1 PM/km: visão telescópica tipo drone. Corpo físico fica Cego e indefeso (CA 10) enquanto em uso. \\
4 & Olho de Permeação & 1 Ação, 2 PM/cena: Raio-X através de paredes/roupas (9m). Não atravessa criaturas vivas ou materiais densos em mana. \\
5 & Olho de Identificação & 1 Ação, 1 PM/alvo: revela fraquezas, nome do feitiço e efeitos. Segredos divinos/de outros continentes aparecem como 'Desconhecido'. \\
6 & Olho da Absorção & Reação, PM igual ao da magia absorvida: anula magia inimiga. Lançar magia com o olho descoberto suga o próprio feitiço (perde ação, PM e a magia falha). \\
7 & Olhos Que Tudo Veem & Ritual de 10min, 10 PM, 1x/semana: rastreia alguém no globo ou revela planta de masmorra. Deixa Visão Embaçada (-2 em acertos físicos) pelo resto do dia. \\
8 & Olho do Vazio Absoluto & 1 Ação, 5 PM/turno mantido: barreira de repulsão de 9m. Não pode atacar ou se mover enquanto mantiver. \\
9 & Olho de Afeição & Passiva se descoberto, 1 PM/hora: quem olha nos seus olhos faz teste de Espírito com Desvantagem ou desenvolve infatuação perigosa (risco de obsessão yandere). \\
10 & Olho Rastreador & 1 Ação de Busca, 2 PM (recente) / 10 PM (décadas): revela rastros de vida, segue pegadas por continentes. Rastreio antigo tem cooldown de 1x/mês. \\
\end{tabelalivro}


% ---------------------------------------------------------------------------
\section{O Valor do Rank e o Bônus de Conjuração}

O quão forte uma magia atinge o inimigo, ou a dificuldade de esquivar de um golpe de
espada, não depende só da força bruta, mas da Maestria (o Rank) naquela escola específica.

\begin{nota}{O Bônus Depende da Ação}
    O Bônus Numérico é específico da árvore em uso no momento. Se você atacar com Magia de
    Água, usa seu Rank de Água; se logo depois usar Magia de Cura, usa o Rank de Cura ---
    que pode ser bem diferente.
\end{nota}

\subsection{O Bônus de Conjuração (BC)}

O número mais importante da ficha de um mago, e ele unifica as três fórmulas do sistema num
único valor.

\begin{formula}
    \textbf{BC $=$ Intelecto $+$ Bônus do Rank naquela escola}
\end{formula}

\begin{itemize}
    \item \textbf{Acerto Mágico:} 1d20 $+$ BC (contra a CA do alvo).
    \item \textbf{CD para resistir:} 8 $+$ BC (o alvo rola 1d20 $+$ atributo de defesa).
    \item \textbf{Dano Mágico:} dados da magia $+$ BC.
\end{itemize}

\subsection{As Fórmulas Marciais}

Guerreiros usam a mesma lógica, trocando o atributo:

\begin{itemize}
    \item \textbf{Acerto Físico} $=$ 1d20 $+$ Força (ou Agilidade, para armas leves) $+$
          Bônus do Rank da Técnica.
    \item \textbf{CD de uma Técnica} $=$ 8 $+$ Força (ou Agilidade) $+$ Bônus do Rank da
          Técnica.
    \item \textbf{Dano Total} $=$ Dados da Arma $+$ Atributo $+$ Bônus do Rank da Técnica.
\end{itemize}

% ---------------------------------------------------------------------------
\section{Misturando Árvores (Multiclasse)}

Este sistema não tem classes --- nada impede você de ser Trovador de Bardo, Rastreador do
Tático, Intermediário de Fogo e Principiante do Norte ao mesmo tempo. É intencional e é o
coração do jogo. Sete perguntas, sete respostas.

\begin{description}
    \item[1. Qual Bônus de Rank eu uso?] O da árvore que concedeu a habilidade, sempre.
          Exceção: quando uma regra genérica pedir seu Bônus de Rank sem dizer de qual
          árvore, use o maior que você possuir em qualquer uma.

    \item[2. Conhecimentos somam entre árvores?] Nunca. A contagem pra desbloquear um
          patamar conta apenas magias, técnicas e talentos daquela mesma árvore.

    \item[3. Como somam PV, PM, PT e PP?] \hfill
          \begin{itemize}
              \item \textbf{PV:} uma reserva só, calculada de uma vez. Os Dados de PV de
                    \emph{todas} as suas árvores entram no mesmo somatório do Cap. 4, §1 ---
                    abrir uma segunda árvore acrescenta dados novos à mesma conta, não uma
                    segunda barra de vida.
              \item \textbf{PM:} somam apenas das escolas de magia, pela fórmula única do
                    Cap. 4.
              \item \textbf{PT:} reserva única, mesmo com vários estilos marciais.
              \item \textbf{PP:} reserva única, mesmo com várias árvores de Utilidade.
          \end{itemize}
          Você nunca tem duas reservas do mesmo tipo.

    \item[4. Custo de Abertura] Abrir uma árvore nova fica mais caro a cada árvore que você
          já tem: \textbf{1ª $=$ 1 PA, 2ª $=$ 2 PA, 3ª $=$ 3 PA, 4ª $=$ 4 PA, 5ª $=$ 5 PA}.
          Cada 1º patamar entrega uma Maestria gratuita --- sem o Custo de Abertura, a
          jogada ótima seria abrir cinco árvores por 5 PA e colecionar cinco Maestrias sem
          nunca subir nenhuma. Agora isso custa 15 PA, e continua sendo uma opção legítima
          --- só não é mais de graça.
\end{description}

\begin{aviso}{5. Largura ou profundidade? (a conta real)}
    \textbf{Profundidade custa 34 PA}, não 12. Os 12 PA são só os desbloqueios
    ($1+1+2+2+3+3$); pra ter \emph{direito} de comprá-los você precisa acumular 10
    conhecimentos naquela mesma árvore, e os 10 mais baratos possíveis custam outros 22 PA
    --- dois de cada patamar, a 1, 1, 2, 3 e 4 PA. Em troca: Bônus de Rank $+6$, seis
    Maestrias, e as magias que só existem lá em cima.

    \medskip
    \textbf{Largura custa 15 PA} por cinco árvores e entrega cinco Maestrias de 1º patamar,
    versatilidade e nenhum teto --- mas trava seu Bônus de Rank em $+1$, o que significa
    errar mais, causar menos dano e ter CDs que qualquer coisa resiste.

    \medskip
    Ou seja: largura é \textbf{mais barata}, profundidade é \textbf{mais forte}. O sistema
    cobra pelos dois caminhos e não deixa nenhum ser simplesmente melhor --- mas é bom saber
    qual dos dois está pedindo mais da sua ficha antes de escolher.
\end{aviso}

\begin{description}
    \item[6. E o Rank Deus?] Sempre narrativo e sempre de uma árvore. Ninguém no Mundo de
          Seis Faces jamais alcançou o patamar Deus em duas coisas ao mesmo tempo --- o
          livro trata isso como impossível, não como difícil.

    \item[7. E se eu for fundo em duas árvores ao mesmo tempo?] Algumas combinações de Rank
          Avançado ou superior revelam uma \termo{árvore híbrida} que não existe pra
          ninguém que não cumpriu os dois pré-requisitos. O Estilo Vendaval (Cap. 3) é a
          primeira: emerge de já dominar o Estilo Deus do Norte e a Magia de Vento. Ela não
          aparece na escolha da Árvore Inicial, e o desbloqueio não é travado por regra
          nenhuma --- o Mestre decide.
\end{description}
